\chapter{Conception du système}
\label{chap:conception}

Ce chapitre traduit la spécification en architecture. Nous présentons d'abord la
structure d'ensemble et les choix qui la fondent, puis la vue statique du système :
composants, déploiement et modèle de données. Vient ensuite la conception des deux
étages d'apprentissage — la chaîne de traitement du signal et les trois architectures
de réseaux, puis le pipeline de prédiction de sévérité. Nous décrivons enfin la vue
dynamique, à travers les enchaînements d'interactions et le cycle de vie d'un examen,
avant d'exposer la conception de la sécurité.

% =============================================================================
\section{Architecture générale}
% =============================================================================

\subsection{Une architecture en trois couches}

\hypnoria{} suit une architecture client-serveur à trois couches, dont la
figure~\ref{fig:archi3t} donne la vue d'ensemble. La couche de présentation est une
application monopage exécutée dans le navigateur du praticien. La couche applicative
expose une interface de programmation sans état, qui porte à la fois la logique métier
et le moteur d'inférence. La couche de persistance sépare délibérément deux natures de
données : les métadonnées cliniques, structurées et interrogeables, vont dans une base
relationnelle ; les fichiers volumineux — enregistrements bruts, images, comptes
rendus — partent vers un stockage objet.

\begin{figure}[H]
\centering
\begin{tikzpicture}[
    x=1cm,y=1cm,
    couche/.style = {rectangle,rounded corners=3pt,draw=hypnoblue!45,
                     fill=hypnolight!60,minimum width=13.6cm},
    mod/.style    = {rectangle,rounded corners=2pt,draw=hypnoblue!65,fill=white,
                     align=center,font=\scriptsize,minimum height=0.95cm,
                     text width=2.45cm},
    modml/.style  = {rectangle,rounded corners=2pt,draw=hypnored!70,fill=hypnored!7,
                     align=center,font=\scriptsize,minimum height=0.95cm,
                     text width=2.45cm},
    ext/.style    = {rectangle,rounded corners=2pt,draw=hypnogray!60,fill=white,
                     align=center,font=\scriptsize,minimum height=0.95cm,
                     text width=2.9cm},
    titre/.style  = {font=\scriptsize\bfseries,text=hypnoblue,anchor=west},
    fl/.style     = {{Latex[length=2mm]}-{Latex[length=2mm]},hypnoblue!70,
                     line width=0.8pt},
    lab/.style    = {font=\tiny\itshape,text=hypnogray,fill=white,inner sep=1.5pt},
  ]

  % ── Couche présentation ─────────────────────────────────────────────────
  \node[couche,minimum height=1.5cm] (c1) at (7.0,4.4) {};
  \node[titre] at (0.25,5.02) {Présentation};
  \node[mod] at (3.3,4.25) {Assistant\\d'analyse};
  \node[mod] at (6.0,4.25) {Registre\\patients};
  \node[mod] at (8.7,4.25) {Console\\d'administration};
  \node[mod] at (11.4,4.25) {Consultations};

  % ── Couche application ──────────────────────────────────────────────────
  \node[couche,minimum height=2.7cm] (c2) at (7.0,1.35) {};
  \node[titre] at (0.25,2.45) {Application};
  \node[mod]   at (3.3,1.75) {Authentification\\et rôles};
  \node[mod]   at (6.0,1.75) {Gestion\\métier};
  \node[mod]   at (8.7,1.75) {Archivage};
  \node[mod]   at (11.4,1.75) {Journalisation};
  \node[modml] at (4.65,0.45) {Prétraitement\\du signal};
  \node[modml] at (7.35,0.45) {Classification\\des stades};
  \node[modml] at (10.05,0.45) {Sévérité\\et explication};
  \node[font=\tiny\itshape,text=hypnored,anchor=east] at (3.30,0.45) {moteur d'inférence};

  % ── Couche persistance ──────────────────────────────────────────────────
  \node[couche,minimum height=1.5cm] (c3) at (7.0,-1.75) {};
  \node[titre] at (0.25,-1.13) {Persistance};
  \node[ext] at (4.6,-1.9) {Base relationnelle\\\tiny métadonnées cliniques};
  \node[ext] at (9.4,-1.9) {Stockage objet\\\tiny fichiers volumineux};

  % ── Flux ────────────────────────────────────────────────────────────────
  \draw[fl] (7.0,3.62) -- node[lab] {interface applicative sur HTTP} (7.0,2.75);
  \draw[fl] (4.6,-0.02) -- node[lab] {requêtes} (4.6,-1.12);
  \draw[fl] (9.4,-0.02) -- node[lab] {liens signés} (9.4,-1.12);

\end{tikzpicture}
\vspace{0.3cm}
\caption[Architecture générale du système]%
        {Architecture en trois couches. Les modules encadrés en rouge constituent le
         moteur d'inférence.}
\label{fig:archi3t}
\end{figure}

\subsection{Justification des choix structurants}

Quatre décisions méritent d'être argumentées, car elles engagent le reste de la
conception. Le tableau~\ref{tab:choix} les récapitule.

\begin{table}[H]
\centering
\caption{Décisions d'architecture et leur justification}
\label{tab:choix}
\small
\begin{tabular}{@{}L{3.2cm}L{5.2cm}L{5.6cm}@{}}
\toprule
\textbf{Décision} & \textbf{Alternative écartée} & \textbf{Motif} \\
\midrule
Application monopage découplée de l'interface applicative &
Application web à rendu serveur &
L'analyse d'un enregistrement dure plusieurs dizaines de secondes ; un client
autonome peut afficher la progression, poursuivre les transferts en arrière-plan et
recomposer les visualisations sans rechargement \\
\addlinespace[3pt]
Moteur d'inférence hébergé dans le même processus que l'interface applicative &
Service d'inférence séparé &
Les modèles pèsent plusieurs centaines de mégaoctets en mémoire ; les charger une
seule fois et les conserver en cache évite un coût de démarrage à chaque requête. La
frontière logique reste posée pour une séparation ultérieure \\
\addlinespace[3pt]
Séparation métadonnées / fichiers volumineux &
Stockage des fichiers en base &
Un enregistrement pèse plusieurs centaines de mégaoctets. Le conserver en base
dégraderait les sauvegardes et les temps de réponse, sans bénéfice transactionnel \\
\addlinespace[3pt]
Authentification par jeton signé, sans session serveur &
Session en mémoire ou en base &
L'absence d'état côté serveur autorise la réplication horizontale sans mécanisme de
partage de session \\
\bottomrule
\end{tabular}
\end{table}

\subsection{Vue de déploiement}

La figure~\ref{fig:deploiement} décrit la topologie d'exécution. Chaque service est
conteneurisé et l'ensemble démarre par une commande unique. Le service de stockage
est le seul élément externe à l'infrastructure.

\begin{figure}[H]
\centering
\begin{tikzpicture}[
    x=1cm,y=1cm,
    noeud/.style  = {rectangle,rounded corners=3pt,draw=hypnoblue!55,fill=white,
                     align=center,font=\scriptsize,minimum height=1.5cm,
                     text width=3.2cm},
    hote/.style   = {rectangle,rounded corners=3pt,draw=hypnogray!55,dashed,
                     fill=hypnolight!35},
    cloud/.style  = {rectangle,rounded corners=8pt,draw=hypnogray!65,fill=white,
                     align=center,font=\scriptsize,minimum height=1.5cm,
                     text width=3.0cm},
    fl/.style     = {-{Latex[length=2mm]},hypnoblue!75,line width=0.8pt},
    lab/.style    = {font=\tiny\itshape,text=hypnogray,fill=white,inner sep=1.5pt},
  ]

  \node[hote,minimum width=11.4cm,minimum height=2.5cm] at (5.7,0) {};
  \node[font=\scriptsize\bfseries,text=hypnogray,anchor=north west] at (0.15,1.15)
       {Hôte de déploiement --- orchestration de conteneurs};

  \node[noeud] (front) at (2.1,-0.2)  {\textbf{Serveur web}\\\tiny Nginx\\\tiny
                                        application compilée};
  \node[noeud] (back)  at (5.7,-0.2)  {\textbf{Interface applicative}\\\tiny Uvicorn
                                        et FastAPI\\\tiny modèles en mémoire};
  \node[noeud] (db)    at (9.3,-0.2)  {\textbf{Base de données}\\\tiny PostgreSQL\\
                                        \tiny volume persistant};

  \node[cloud] (b2)    at (13.6,-0.2) {\textbf{Stockage objet}\\\tiny service externe\\
                                        \tiny interface S3};

  \node[font=\scriptsize\bfseries,align=center,text width=2.2cm] (nav) at (2.1,2.35)
       {Navigateur\\du praticien};

  \draw[fl] (nav) -- node[lab] {HTTPS} (front);
  \draw[fl] (front) -- node[lab] {mandataire} (back);
  \draw[fl] (back) -- node[lab] {SQL} (db);
  \draw[fl] (back) -- node[lab] {S3} (b2);

\end{tikzpicture}
\vspace{0.3cm}
\caption[Vue de déploiement]%
        {Topologie d'exécution : trois conteneurs sur un hôte, un service de stockage
         externe.}
\label{fig:deploiement}
\end{figure}

% =============================================================================
\section{Conception de la base de données}
% =============================================================================

\subsection{Diagramme de classes}

Le modèle métier s'organise autour de huit entités
(figure~\ref{fig:classes}). L'entité centrale est l'\emph{examen}, qui matérialise une
nuit d'enregistrement : elle porte les références vers les fichiers archivés, la
classe de sévérité prédite et le document de résultats. Autour d'elle gravitent le
patient qui l'a subie, les annotations que le praticien y a portées et les
discussions qu'elle a suscitées.

\begin{figure}[H]
\centering
\begin{tikzpicture}[
    x=1cm,y=1cm,
    cls/.style  = {rectangle,draw=hypnoblue!65,fill=white,align=left,anchor=north,
                   font=\tiny,inner sep=0pt,text width=3.5cm},
    ent/.style  = {rectangle,draw=hypnored!70,fill=hypnored!6,align=left,anchor=north,
                   font=\tiny,inner sep=0pt,text width=3.5cm},
    rel/.style  = {hypnogray!85,line width=0.6pt},
    card/.style = {font=\tiny,text=hypnogray,inner sep=1.6pt,fill=white},
    note/.style = {font=\tiny\itshape,text=hypnogray,align=center,text width=3.5cm},
  ]

  \newcommand{\classe}[3]{%
    \begin{tabular}{@{\hspace{3pt}}p{3.3cm}@{}}
      \rowcolor{hypnolight}\textbf{#1}\\[1pt]
      \hline
      \rule{0pt}{7pt}#2\\[1pt]
      \hline
      \rule{0pt}{7pt}\textit{#3}\\[1pt]
    \end{tabular}}

  % ── Rangée 1 ────────────────────────────────────────────────────────────
  \node[cls] (hop) at (2.0, 4.40) {\classe{Hospital}{id, name, billing\_tier,
       b2\_bucket, created\_at}{--}};
  \node[cls] (usr) at (7.0, 4.40) {\classe{User}{id, email, hashed\_password, role,
       first\_name, last\_name, status, license\_expiry, last\_login}{login(),
       logout()}};
  \node[cls] (pat) at (12.0, 4.40) {\classe{Patient}{id, first\_name, last\_name,
       age, imc, gender}{--}};

  % ── Rangée 2 ────────────────────────────────────────────────────────────
  \node[cls] (log)  at (2.0, 0.55) {\classe{AuditLog}{id, timestamp, user\_email,
       user\_role, action, ip\_address}{--}};
  \node[cls] (conv) at (7.0, 0.55) {\classe{FileConversation}{id, file\_type,
       created\_at}{--}};
  \node[ent] (psg)  at (12.0, 0.55) {\classe{PSG}{id, date, severity, report\_data,
       edf\_url, hypnogram\_url, hypnogram\_annotated\_url, csv\_url,
       osa\_report\_url}{--}};

  % ── Rangée 3 ────────────────────────────────────────────────────────────
  \node[cls] (msg) at (7.0, -3.35) {\classe{FileMessage}{id, content, timestamp}{--}};
  \node[cls] (ann) at (12.0,-3.35) {\classe{PSGAnnotation}{id, epoch\_index, note,
       created\_at}{--}};

  \node[note,anchor=north] at (2.0,-1.05)
       {table volontairement dénormalisée :\\aucune clé étrangère, afin que la
        trace survive à la suppression du compte};

  % ── Associations ────────────────────────────────────────────────────────
  \draw[rel] (hop.east) -- node[card,pos=0.12] {1}
                           node[card,pos=0.88] {0..*} (usr.west);
  \draw[rel] (usr.east) -- node[card,pos=0.12] {1}
                           node[card,pos=0.88] {0..*} (pat.west);

  \draw[rel] (pat.south) -- node[card,pos=0.13] {1}
                            node[card,pos=0.87] {0..*} (psg.north);
  \draw[rel] (psg.south) -- node[card,pos=0.13] {1}
                            node[card,pos=0.87] {0..*} (ann.north);

  \draw[rel] (psg.west) -- node[card,pos=0.13] {1}
                           node[card,pos=0.87] {0..*} (conv.east);
  \draw[rel] (conv.south) -- node[card,pos=0.13] {1}
                             node[card,pos=0.87] {0..*} (msg.north);

  \draw[rel] (usr.south) -- node[card,pos=0.13] {2}
                            node[card,pos=0.88] {0..*} (conv.north);

  % Auteur d'un message : contournement par la gauche
  \draw[rel] (usr.west) -- ++(-0.55,0)
             node[card,pos=0.55] {1}
             |- node[card,pos=0.93] {0..*} (msg.west);

\end{tikzpicture}
\vspace{0.2cm}
\caption[Diagramme de classes du modèle métier]%
        {Modèle métier : l'examen est l'entité pivot du système.}
\label{fig:classes}
\end{figure}

\subsection{Dictionnaire de données}

Le tableau~\ref{tab:dico} décrit les tables et leurs particularités. Le dictionnaire
complet, colonne par colonne, figure en annexe.

\begin{table}[H]
\centering
\caption{Tables du schéma relationnel}
\label{tab:dico}
\small
\begin{tabular}{@{}lcL{8.5cm}@{}}
\toprule
\textbf{Table} & \textbf{Clés étrangères} & \textbf{Rôle et particularités} \\
\midrule
\texttt{hospitals} & --- &
Établissements de rattachement ; prévoit un compartiment de stockage propre à chaque
structure \\
\addlinespace[2pt]
\texttt{users} & \texttt{hospital\_id} &
Comptes unifiés pour les deux rôles ; porte le statut et l'échéance de licence,
tous deux évalués à la connexion \\
\addlinespace[2pt]
\texttt{patients} & \texttt{doctor\_id} &
Rattachement direct au praticien, qui fonde le cloisonnement des données \\
\addlinespace[2pt]
\texttt{psgs} & \texttt{patient\_id} &
Entité pivot ; ne stocke que des références vers les fichiers, jamais leur contenu \\
\addlinespace[2pt]
\texttt{psg\_annotations} & \texttt{psg\_id} &
Observations horodatées à l'époque ; suppression en cascade avec l'examen \\
\addlinespace[2pt]
\texttt{audit\_logs} & --- &
Journal volontairement dénormalisé : l'adresse électronique et le rôle sont recopiés,
afin que la trace survive à la suppression du compte \\
\addlinespace[2pt]
\texttt{file\_conversations} & \texttt{psg\_id}, deux \texttt{user\_id} &
Fil de discussion rattaché à un examen et à exactement deux praticiens \\
\addlinespace[2pt]
\texttt{file\_messages} & \texttt{conversation\_id}, \texttt{sender\_id} &
Messages horodatés ; suppression en cascade avec la discussion \\
\bottomrule
\end{tabular}
\end{table}

Deux choix méritent un commentaire. Le \textbf{journal d'audit est dénormalisé} : il
recopie l'adresse et le rôle de l'auteur au lieu de référencer son compte. C'est
délibéré — une trace de conformité doit rester lisible même après la suppression du
compte concerné, et une clé étrangère l'en empêcherait.

Les \textbf{références de fichiers sont stockées en clair} dans la table des examens,
mais ne sont jamais transmises telles quelles : elles sont converties en liens signés
à durée limitée au moment de la sérialisation de la réponse. Ce point est développé à
la section~\ref{sec:securite}.

% =============================================================================
\section{Conception de la chaîne de classification}
% =============================================================================

\subsection{Traitement du signal}

Avant toute inférence, le signal brut doit être ramené à une forme normalisée. La
figure~\ref{fig:chaine} décrit cette chaîne, dont chaque étape répond à une source de
variabilité identifiée.

\begin{figure}[H]
\centering
\begin{tikzpicture}[
    x=1cm,y=1cm,
    et/.style  = {rectangle,rounded corners=2pt,draw=hypnoblue!60,fill=hypnolight,
                  align=center,font=\scriptsize,minimum height=1.25cm,
                  text width=2.15cm},
    fl/.style  = {-{Latex[length=2mm]},hypnoblue!75,line width=0.8pt},
    note/.style= {font=\tiny\itshape,text=hypnogray,align=center,text width=2.3cm},
  ]

  \node[et] (a) at (0.0,0)  {Décodage\\du fichier};
  \node[et] (b) at (2.75,0) {Résolution\\des dérivations};
  \node[et] (c) at (5.5,0)  {Rééchantillonnage\\à 100\,Hz};
  \node[et] (d) at (8.25,0) {Segmentation\\en époques};
  \node[et] (e) at (11.0,0) {Normalisation\\par époque};
  \node[et] (f) at (13.75,0){Tenseur\\$(N, C, 3000)$};

  \foreach \x/\y in {a/b,b/c,c/d,d/e,e/f} \draw[fl] (\x) -- (\y);

  \node[note,anchor=north] at (0.0,-0.78)  {format\\constructeur};
  \node[note,anchor=north] at (2.75,-0.78) {conventions de\\nommage};
  \node[note,anchor=north] at (5.5,-0.78)  {fréquences\\d'acquisition};
  \node[note,anchor=north] at (8.25,-0.78) {convention\\AASM 30\,s};
  \node[note,anchor=north] at (11.0,-0.78) {amplitudes\\inter-sujets};
  \node[note,anchor=north] at (13.75,-0.78){entrée des\\modèles};

  \node[font=\tiny\bfseries,text=hypnored,anchor=east] at (-1.05,-1.35)
       {variabilité traitée :};

\end{tikzpicture}
\vspace{0.3cm}
\caption[Chaîne de traitement du signal]%
        {Chaîne de préparation : chaque étape neutralise une source de variabilité.}
\label{fig:chaine}
\end{figure}

L'étape de \textbf{normalisation} appelle une justification particulière, car elle
conditionne la validité de l'évaluation. La statistique de centrage et de réduction
est calculée sur chaque segment isolément, et non sur le corpus. Ce choix neutralise
les écarts d'amplitude liés à l'impédance des électrodes, au gain de l'amplificateur
et à la morphologie du sujet, ainsi que la dérive lente au cours de la nuit. Il a
surtout une conséquence méthodologique : \textbf{aucune statistique globale n'étant
estimée, aucune information ne peut fuir de l'ensemble d'entraînement vers l'ensemble
de test par la normalisation}. En contrepartie, l'information d'amplitude absolue est
perdue, alors qu'elle caractérise partiellement le sommeil profond ; c'est la forme et
la fréquence des oscillations qui devront porter cette discrimination.

\subsection{Architecture convolutive}

Le premier modèle de base exploite l'observation formulée au
chapitre~\ref{chap:etat-art} : les éléments graphiques du référentiel AASM vivent à
des échelles temporelles différentes. Un fuseau du sommeil dure environ une demi-seconde
et oscille autour de 12 à 14 hertz ; une onde lente s'étale sur plus d'une seconde.
Un noyau de convolution unique ne peut capter les deux.

L'architecture retenue (figure~\ref{fig:cnn}) comporte donc deux branches parallèles
appliquées au même signal. La branche fine emploie un noyau de 50 échantillons — une
demi-seconde — et la branche grossière un noyau de 200 échantillons, soit deux
secondes. Leurs représentations sont concaténées avant classification.

\begin{figure}[H]
\centering
\begin{tikzpicture}[
    x=1cm,y=1cm,
    bl/.style   = {rectangle,rounded corners=2pt,draw=hypnoblue!60,fill=hypnolight,
                   align=center,font=\tiny,minimum height=0.95cm,text width=2.0cm},
    blf/.style  = {rectangle,rounded corners=2pt,draw=hypnored!65,fill=hypnored!7,
                   align=center,font=\tiny,minimum height=0.95cm,text width=2.0cm},
    io/.style   = {rectangle,rounded corners=2pt,draw=hypnogray!60,fill=white,
                   align=center,font=\tiny,minimum height=0.95cm,text width=1.9cm},
    fl/.style   = {-{Latex[length=1.8mm]},hypnoblue!75,line width=0.7pt},
    tag/.style  = {font=\tiny\itshape,text=hypnogray},
  ]

  \node[io] (in) at (0,0) {Entrée\\$(C, 3000)$};

  \node[bl] (f1) at (3.0, 1.35) {Conv 1D\\noyau 50, pas 6};
  \node[bl] (f2) at (5.6, 1.35) {Sous-échant.\\$\times 4$};
  \node[bl] (f3) at (8.2, 1.35) {Conv 1D\\noyau 8};
  \node[bl] (f4) at (10.8,1.35) {Moyennage\\adaptatif $\to 16$};

  \node[bl] (g1) at (3.0,-1.35) {Conv 1D\\noyau 200, pas 25};
  \node[bl] (g2) at (5.6,-1.35) {Sous-échant.\\$\times 4$};
  \node[bl] (g3) at (8.2,-1.35) {Conv 1D\\noyau 8};
  \node[bl] (g4) at (10.8,-1.35) {Moyennage\\adaptatif $\to 16$};

  \node[blf] (cat) at (13.3,0) {Concaténation\\$4096$};
  \node[io]  (out) at (13.3,-2.1) {Logits\\$(K)$};

  \foreach \a/\b in {f1/f2,f2/f3,f3/f4,g1/g2,g2/g3,g3/g4} \draw[fl] (\a) -- (\b);
  \draw[fl] (in) |- (f1);
  \draw[fl] (in) |- (g1);
  \draw[fl] (f4) -| (cat);
  \draw[fl] (g4) -| (cat);
  \draw[fl] (cat) -- node[tag,right,pos=0.5] {\ 3 couches denses} (out);

  \node[tag,anchor=west] at (1.35, 1.95) {branche fine --- 0,5\,s};
  \node[tag,anchor=west] at (1.35,-1.95) {branche grossière --- 2\,s};

\end{tikzpicture}
\vspace{0.3cm}
\caption[Architecture du réseau convolutif]%
        {Réseau convolutif à deux branches d'échelles temporelles distinctes.}
\label{fig:cnn}
\end{figure}

\subsection{Architecture récurrente à pooling attentionnel}

Le deuxième modèle traite l'époque comme une séquence de 3\,000 pas, chacun décrit par
les $C$ dérivations. Un réseau récurrent bidirectionnel à mémoire longue, de deux
couches de 128 unités, en produit une représentation contextualisée dans les deux sens
du temps.

Se pose alors la question de la réduction : comment condenser 3\,000 états en un
vecteur unique ? Prendre le dernier état privilégierait la fin de la fenêtre ; en
faire la moyenne diluerait un événement bref dans le reste. Nous employons un
\textbf{pooling attentionnel} : un petit réseau apprend à attribuer à chaque pas un
poids, dont la somme vaut un, et la représentation finale est la moyenne pondérée des
états (figure~\ref{fig:lstm}). Le modèle apprend ainsi \emph{où regarder}, ce qui
convient à un signal où l'information discriminante est localisée.

\begin{figure}[H]
\centering
\begin{tikzpicture}[
    x=1cm,y=1cm,
    bl/.style   = {rectangle,rounded corners=2pt,draw=hypnoblue!60,fill=hypnolight,
                   align=center,font=\tiny,minimum height=1.0cm,text width=2.3cm},
    blf/.style  = {rectangle,rounded corners=2pt,draw=hypnored!65,fill=hypnored!7,
                   align=center,font=\tiny,minimum height=1.0cm,text width=2.3cm},
    io/.style   = {rectangle,rounded corners=2pt,draw=hypnogray!60,fill=white,
                   align=center,font=\tiny,minimum height=1.0cm,text width=2.0cm},
    fl/.style   = {-{Latex[length=1.8mm]},hypnoblue!75,line width=0.7pt},
    tag/.style  = {font=\tiny\itshape,text=hypnogray,align=center},
  ]

  \node[io] (in)  at (0,0)    {Séquence\\$(3000, C)$};
  \node[bl] (l1)  at (3.1,0)  {LSTM bidirectionnel\\2 couches $\times$ 128};
  \node[bl] (st)  at (6.4,0)  {États contextualisés\\$(3000, 256)$};
  \node[blf](at)  at (9.7,0)  {Pooling attentionnel\\$\sum_t \alpha_t h_t$};
  \node[bl] (fc)  at (12.9,0) {Couches denses\\$256 \to 64 \to K$};

  \foreach \a/\b in {in/l1,l1/st,st/at,at/fc} \draw[fl] (\a) -- (\b);

  \node[bl,minimum height=0.75cm,text width=2.3cm] (sc) at (9.7,-1.9)
       {Scores $\alpha_t$\\\tiny normalisés, de somme 1};
  \draw[fl] (st) |- (sc);
  \draw[fl] (sc) -- (at);

  \node[tag,anchor=north,text width=6.5cm] at (6.4,-2.35)
       {les poids d'attention indiquent les instants retenus};

\end{tikzpicture}
\vspace{0.3cm}
\caption[Architecture du réseau récurrent]%
        {Réseau récurrent bidirectionnel à pooling attentionnel.}
\label{fig:lstm}
\end{figure}

\subsection{Architecture attentionnelle}

Le troisième modèle abandonne la récurrence. L'époque est découpée en 60 fragments
consécutifs de 50 échantillons, chacun projeté linéairement dans un espace de
dimension 128. Un jeton de synthèse est ajouté en tête de séquence, et un encodage
positionnel sinusoïdal réinjecte l'ordre temporel que l'auto-attention, par
construction, ignore. Quatre couches d'encodeur à huit têtes traitent l'ensemble ; la
classification s'opère sur le seul jeton de synthèse
(figure~\ref{fig:transformer}).

Le découpage en fragments est ce qui rend l'approche viable. L'auto-attention a un
coût quadratique en la longueur de séquence : l'appliquer aux 3\,000 échantillons bruts
demanderait neuf millions de coefficients par tête et par époque. Ramenée à 61 jetons,
la même opération en demande moins de quatre mille.

\begin{figure}[H]
\centering
\begin{tikzpicture}[
    x=1cm,y=1cm,
    bl/.style   = {rectangle,rounded corners=2pt,draw=hypnoblue!60,fill=hypnolight,
                   align=center,font=\tiny,minimum height=1.0cm,text width=2.25cm},
    blf/.style  = {rectangle,rounded corners=2pt,draw=hypnored!65,fill=hypnored!7,
                   align=center,font=\tiny,minimum height=1.0cm,text width=2.25cm},
    io/.style   = {rectangle,rounded corners=2pt,draw=hypnogray!60,fill=white,
                   align=center,font=\tiny,minimum height=1.0cm,text width=1.95cm},
    fl/.style   = {-{Latex[length=1.8mm]},hypnoblue!75,line width=0.7pt},
    tag/.style  = {font=\tiny\itshape,text=hypnogray,align=center},
  ]

  \node[io]  (in)  at (0,0)    {Entrée\\$(C, 3000)$};
  \node[bl]  (pa)  at (2.9,0)  {Découpage\\60 fragments de 50};
  \node[bl]  (pe)  at (5.9,0)  {Projection linéaire\\$\to 128$};
  \node[blf] (cls) at (8.9,0)  {Jeton de synthèse\\et encodage positionnel};
  \node[bl]  (tr)  at (11.9,0) {Encodeur\\4 couches, 8 têtes};
  \node[io]  (out) at (11.9,-2.0) {Logits\\$(K)$};

  \foreach \a/\b in {in/pa,pa/pe,pe/cls,cls/tr} \draw[fl] (\a) -- (\b);
  \draw[fl] (tr) -- node[tag,right,pos=0.5,text width=2.4cm]
                    {\ classification sur le jeton de synthèse} (out);

  \node[tag,anchor=north,text width=5.0cm] at (4.4,-1.05)
       {coût quadratique ramené de $3000^2$ à $61^2$};

\end{tikzpicture}
\vspace{0.3cm}
\caption[Architecture du modèle attentionnel]%
        {Modèle attentionnel : l'époque est traitée comme une suite de fragments.}
\label{fig:transformer}
\end{figure}

\subsection{Combinaison par empilement}

Les trois modèles sont ensuite combinés selon le principe exposé à la
figure~\ref{fig:stacking-principe}. Chacun produit un vecteur de probabilités sur les
$K$ classes ; les trois vecteurs sont concaténés en $3K$ méta-caractéristiques qui
alimentent une régression logistique équilibrée. Le choix d'un modèle linéaire pour le
second niveau n'est pas un défaut d'ambition : le méta-apprenant dispose de peu de
variables et doit avant tout produire des probabilités calibrées, ce qu'un modèle
simple fait mieux qu'un réseau supplémentaire.

Quatre ensembles sont construits, un par configuration de montage et de granularité.

% =============================================================================
\section{Conception du pipeline de sévérité}
% =============================================================================

Le second étage part de l'hypnogramme produit et non du signal. La
figure~\ref{fig:pipeline-osa} en décrit l'enchaînement.

\begin{figure}[H]
\centering
\begin{tikzpicture}[
    x=1cm,y=1cm,
    et/.style   = {rectangle,rounded corners=2pt,draw=hypnoblue!60,fill=hypnolight,
                   align=center,font=\tiny,minimum height=1.15cm,text width=2.3cm},
    etf/.style  = {rectangle,rounded corners=2pt,draw=hypnored!65,fill=hypnored!7,
                   align=center,font=\tiny,minimum height=1.15cm,text width=2.3cm},
    io/.style   = {rectangle,rounded corners=2pt,draw=hypnogray!60,fill=white,
                   align=center,font=\tiny,minimum height=1.15cm,text width=2.1cm},
    fl/.style   = {-{Latex[length=1.8mm]},hypnoblue!75,line width=0.7pt},
    tag/.style  = {font=\tiny\itshape,text=hypnogray,align=center},
  ]

  \node[io] (hyp) at (0,0.8)   {Hypnogramme};
  \node[io] (cli) at (0,-0.9)  {Données\\cliniques};

  \node[et] (mrq) at (3.1,0.8)  {25 marqueurs\\d'architecture};
  \node[et] (fus) at (6.2,0)    {Fusion et\\8 interactions};
  \node[et] (imp) at (9.3,0)    {Imputation par\\la médiane};
  \node[etf](mod) at (12.4,0)   {Ensemble\\par empilement};

  \node[etf](shp) at (12.4,-2.3) {Attribution des\\contributions};
  \node[io] (res) at (8.6,-2.3)  {Classe, probabilités\\et facteurs};

  \draw[fl] (hyp) -- (mrq);
  \draw[fl] (mrq) -- (fus);
  \draw[fl] (cli) -| (fus);
  \draw[fl] (fus) -- (imp);
  \draw[fl] (imp) -- (mod);
  \draw[fl] (mod) -- (shp);
  \draw[fl] (shp) -- (res);

  \node[tag,anchor=north,text width=3.6cm] at (6.2,-1.05)
       {53 variables au total};

\end{tikzpicture}
\vspace{0.3cm}
\caption[Pipeline de prédiction de la sévérité]%
        {Chaîne de prédiction de la sévérité, de l'hypnogramme à l'explication.}
\label{fig:pipeline-osa}
\end{figure}

Trois décisions de conception structurent cette chaîne.

La première concerne le \textbf{choix des variables}. Les index de perturbation
respiratoire disponibles dans la cohorte ont été délibérément exclus : ils mesurent
directement les événements que l'on cherche à prédire, et les inclure aurait produit
un modèle excellent mais sans valeur. Ce point est développé au
chapitre~\ref{chap:realisation-ia}.

La deuxième concerne les \textbf{valeurs manquantes}. Toutes les variables cliniques
ne sont pas systématiquement disponibles en pratique. Plutôt que de refuser la
prédiction, le système comble les absences par la médiane observée sur la population
d'apprentissage, et signale au praticien quelles variables ont été imputées.

La troisième concerne l'\textbf{explication}. La prédiction est rendue par un ensemble
hétérogène, mais l'explication est calculée sur un modèle à arbres entraîné sur les
mêmes données. Ce dédoublement est assumé : il permet d'obtenir des attributions
\emph{exactes} en temps polynomial, là où expliquer directement l'ensemble aurait
imposé une approximation par échantillonnage. Sa limite doit être énoncée dans le
compte rendu — l'explication décrit un modèle très proche du modèle décisionnel, non
ce modèle lui-même.

% =============================================================================
\section{Vue dynamique}
% =============================================================================

\subsection{Authentification}

La figure~\ref{fig:seq-auth} décrit l'ouverture de session. Trois vérifications
successives conditionnent la délivrance du jeton : la validité du mot de passe, le
statut du compte et l'échéance de la licence.

\begin{figure}[H]
\centering
\resizebox{0.98\textwidth}{!}{%
\begin{sequencediagram}
  \newthread[hypnolight]{u}{\scriptsize Praticien}
  \newinst[2.2]{f}{\scriptsize Client web}
  \newinst[2.2]{a}{\scriptsize Interface applicative}
  \newinst[2.2]{d}{\scriptsize Base de données}

  \begin{call}{u}{\scriptsize saisit ses identifiants}{f}{}
  \end{call}
  \begin{call}{f}{\scriptsize demande de jeton}{a}{\scriptsize jeton et profil}
    \begin{call}{a}{\scriptsize recherche du compte}{d}{\scriptsize empreinte}
    \end{call}
    \begin{sdblock}{\scriptsize Vérifications}{\scriptsize mot de passe, statut, licence}
      \begin{call}{a}{\scriptsize horodate la connexion}{d}{}
      \end{call}
      \begin{call}{a}{\scriptsize journalise l'accès}{d}{}
      \end{call}
    \end{sdblock}
  \end{call}
  \begin{call}{f}{\scriptsize conserve le jeton}{f}{}
  \end{call}
  \begin{call}{f}{\scriptsize ouvre l'espace de travail}{u}{}
  \end{call}
\end{sequencediagram}}
\caption[Séquence d'authentification]%
        {Ouverture de session et délivrance du jeton.}
\label{fig:seq-auth}
\end{figure}

\subsection{Analyse d'un enregistrement}

L'enchaînement le plus caractéristique du système est celui de l'analyse
(figure~\ref{fig:seq-analyse}). Sa particularité tient à l'ordre des opérations :
les résultats sont affichés au praticien \emph{avant} que l'archivage du fichier ne
soit terminé. L'enregistrement pèse plusieurs centaines de mégaoctets et son transfert
vers le stockage prend un temps que le médecin n'a aucune raison d'attendre.

\begin{figure}[H]
\centering
\resizebox{0.98\textwidth}{!}{%
\begin{sequencediagram}
  \newthread[hypnolight]{u}{\scriptsize Praticien}
  \newinst[1.8]{f}{\scriptsize Client web}
  \newinst[1.8]{a}{\scriptsize Interface applicative}
  \newinst[1.8]{m}{\scriptsize Moteur d'inférence}
  \newinst[1.8]{d}{\scriptsize Base de données}
  \newinst[1.8]{s}{\scriptsize Stockage objet}

  \begin{call}{u}{\scriptsize dépose le fichier}{f}{}
  \end{call}
  \begin{call}{f}{\scriptsize demande d'analyse}{a}{\scriptsize stades et métriques}
    \begin{call}{a}{\scriptsize prétraitement}{m}{\scriptsize tenseur d'époques}
    \end{call}
    \begin{call}{a}{\scriptsize classification}{m}{\scriptsize séquence de stades}
    \end{call}
    \begin{call}{a}{\scriptsize métriques AASM}{m}{\scriptsize indicateurs}
    \end{call}
  \end{call}
  \begin{call}{f}{\scriptsize affiche l'hypnogramme}{u}{}
  \end{call}
  \begin{call}{f}{\scriptsize enregistre l'examen}{a}{\scriptsize identifiant}
    \begin{call}{a}{\scriptsize insertion}{d}{}
    \end{call}
  \end{call}
  \begin{sdblock}{\scriptsize Arrière-plan}{\scriptsize sans blocage de l'interface}
    \begin{call}{f}{\scriptsize transfert du fichier}{a}{\scriptsize référence}
      \begin{call}{a}{\scriptsize dépôt}{s}{\scriptsize clé d'objet}
      \end{call}
      \begin{call}{a}{\scriptsize mise à jour}{d}{}
      \end{call}
    \end{call}
  \end{sdblock}
\end{sequencediagram}}
\caption[Séquence d'analyse d'un enregistrement]%
        {Analyse d'un enregistrement : les résultats précèdent l'archivage.}
\label{fig:seq-analyse}
\end{figure}

\subsection{Prédiction de la sévérité}

La figure~\ref{fig:seq-osa} décrit la seconde phase, déclenchée depuis les résultats
d'analyse.

\begin{figure}[H]
\centering
\resizebox{0.93\textwidth}{!}{%
\begin{sequencediagram}
  \newthread[hypnolight]{u}{\scriptsize Praticien}
  \newinst[2.2]{f}{\scriptsize Client web}
  \newinst[2.2]{a}{\scriptsize Interface applicative}
  \newinst[2.2]{m}{\scriptsize Moteur d'inférence}

  \begin{call}{u}{\scriptsize saisit les données cliniques}{f}{}
  \end{call}
  \begin{call}{f}{\scriptsize demande d'évaluation}{a}{\scriptsize sévérité et facteurs}
    \begin{call}{a}{\scriptsize extraction des marqueurs}{a}{}
    \end{call}
    \begin{call}{a}{\scriptsize imputation et alignement}{a}{}
    \end{call}
    \begin{call}{a}{\scriptsize prédiction}{m}{\scriptsize classe et probabilités}
    \end{call}
    \begin{call}{a}{\scriptsize attribution des contributions}{m}{\scriptsize facteurs}
    \end{call}
  \end{call}
  \begin{call}{f}{\scriptsize affiche le rapport expliqué}{u}{}
  \end{call}
\end{sequencediagram}}
\caption[Séquence de prédiction de la sévérité]%
        {Prédiction de la sévérité et production des facteurs explicatifs.}
\label{fig:seq-osa}
\end{figure}

\subsection{Cycle de vie d'un examen}

Un examen traverse une succession d'états, décrits à la
figure~\ref{fig:etats}. Cette modélisation explicite un point important : un examen
reste consultable même si l'archivage échoue, car les résultats d'analyse sont
stockés en base indépendamment des fichiers.

\begin{figure}[H]
\centering
\begin{tikzpicture}[
    x=1cm,y=1cm,
    st/.style   = {rectangle,rounded corners=5pt,draw=hypnoblue!65,fill=hypnolight,
                   align=center,font=\scriptsize,minimum height=0.9cm,text width=2.3cm},
    stf/.style  = {rectangle,rounded corners=5pt,draw=hypnored!70,fill=hypnored!7,
                   align=center,font=\scriptsize,minimum height=0.9cm,text width=2.3cm},
    fl/.style   = {-{Latex[length=2mm]},hypnoblue!75,line width=0.7pt},
    tr/.style   = {font=\tiny\itshape,text=hypnogray,fill=white,inner sep=1.5pt},
  ]

  \fill[hypnoblue] (0,0) circle (0.14);

  \node[st]  (cre) at (2.5,0)   {Créé};
  \node[st]  (ana) at (6.0,0)   {Analysé};
  \node[st]  (arc) at (9.5,0)   {Archivé};
  \node[stf] (eva) at (13.0,0)  {Évalué};
  \node[st]  (ann) at (9.5,-2.3){Annoté};
  \node[stf] (deg) at (6.0,-2.3){Archivage\\en échec};

  \draw[fl] (0.16,0) -- node[tr,above] {dépôt du fichier} (cre);
  \draw[fl] (cre) -- node[tr,above] {classification} (ana);
  \draw[fl] (ana) -- node[tr,above] {transfert réussi} (arc);
  \draw[fl] (arc) -- node[tr,above] {prédiction} (eva);
  \draw[fl] (arc) -- node[tr,right,pos=0.5] {annotation} (ann);
  \draw[fl] (ana) -- node[tr,right,pos=0.5] {transfert échoué} (deg);
  \draw[fl] (deg) -- node[tr,below] {nouvelle tentative} (ann);

  \fill[hypnoblue] (13.0,-2.3) circle (0.19);
  \fill[white]     (13.0,-2.3) circle (0.135);
  \fill[hypnoblue] (13.0,-2.3) circle (0.09);
  \draw[fl] (eva) -- node[tr,right,pos=0.5] {compte rendu} (13.0,-2.11);
  \draw[fl] (ann) -- node[tr,pos=0.72,below right] {prédiction} (eva.south west);

\end{tikzpicture}
\vspace{0.3cm}
\caption[Cycle de vie d'un examen]%
        {États successifs d'un examen ; l'échec d'archivage ne bloque pas
         l'exploitation des résultats.}
\label{fig:etats}
\end{figure}

% =============================================================================
\section{Conception de la sécurité}
\label{sec:securite}
% =============================================================================

La sécurité repose sur quatre mécanismes complémentaires, résumés au
tableau~\ref{tab:secu}.

\begin{table}[H]
\centering
\caption{Mécanismes de sécurité et menaces couvertes}
\label{tab:secu}
\small
\begin{tabular}{@{}L{3.2cm}L{5.3cm}L{5.5cm}@{}}
\toprule
\textbf{Mécanisme} & \textbf{Principe} & \textbf{Menace couverte} \\
\midrule
Empreintes de mots de passe &
Fonction de dérivation lente avec sel aléatoire par compte &
Compromission de la base : les mots de passe restent inexploitables \\
\addlinespace[3pt]
Jetons signés &
Jeton porteur signé, transmis à chaque appel, contenant l'identité et le rôle &
Usurpation : toute altération invalide la signature \\
\addlinespace[3pt]
Contrôle d'accès par rôle &
Vérification du rôle et de la propriété de la ressource sur chaque route protégée &
Élévation de privilège et accès aux patients d'un confrère \\
\addlinespace[3pt]
Liens d'accès signés &
Les références de fichiers sont converties en URL signées valables 24 heures &
Fuite de données : aucun objet n'est accessible sans signature valide \\
\bottomrule
\end{tabular}
\end{table}

Le \textbf{cycle de vie des licences} constitue un cinquième niveau, propre au contexte
clinique. Un compte peut être actif, suspendu ou en attente, et porte une date
d'échéance. Ces trois conditions sont évaluées à chaque ouverture de session, ce qui
permet de révoquer un accès sans supprimer le compte ni ses données.

Le \textbf{journal d'audit} complète le dispositif par la traçabilité. Chaque action
sensible y inscrit son auteur, son horodatage, sa nature et l'adresse d'origine. Le
choix de dénormalisation évoqué plus haut prend ici tout son sens : la trace doit
survivre à la disparition de ce qu'elle référence.

% =============================================================================
\section{Conclusion}
% =============================================================================

Ce chapitre a défini l'architecture du système et les structures qui la composent.
L'organisation en trois couches sépare la présentation, la logique métier et la
persistance, cette dernière distinguant elle-même les métadonnées interrogeables des
fichiers volumineux. Le modèle de données s'articule autour de l'examen, entité pivot
qui ne stocke que des références vers les fichiers archivés.

La conception des deux étages d'apprentissage a été guidée par des considérations
explicites : les échelles temporelles des éléments du référentiel AASM pour
l'architecture convolutive, la localisation de l'information discriminante pour le
pooling attentionnel, le coût quadratique de l'auto-attention pour le découpage en
fragments, et la prévention des fuites de la variable cible pour le pipeline de
sévérité.

La vue dynamique a mis en évidence un choix d'enchaînement caractéristique — afficher
les résultats avant la fin de l'archivage — et un principe de robustesse : l'échec
d'un transfert ne compromet pas l'exploitation d'un examen.

Le chapitre suivant passe à la réalisation de ces modèles : préparation des données,
stratégie d'entraînement sous contrainte matérielle et résultats expérimentaux.
